\chapter{Local RX, regularisation and MNF-RX}\label{ch:localrx}

\begin{onething}
Global RX fails in two ways: one Gaussian cannot describe a scene of many surfaces, and a 165-band
covariance cannot be inverted safely. Local RX fixes the first by comparing each pixel with a ring of
its own neighbours; MNF-RX fixes the second by keeping only the few directions in which signal beats
noise, so the covariance it inverts has no small eigenvalues to explode. Both are the Mahalanobis
distance of \cref{ch:rx}, applied to a better-chosen background or a better-chosen space.
\end{onething}

\section{Two failures, two repairs}\label{sec:localrx:motivation}

\Cref{ch:rx} ended with global RX failing for two separate reasons.

\emph{The background is not one cloud.} In \cref{we:rx:mixture} a hot patch of machinery on warm bare
ground scored $z = 1.57$ against a scene that was half forest, and $z = 3.0$ against the bare ground
around it. Global RX asks ``is this pixel unusual for the scene?''; the question you often mean is
``is this pixel unusual for its surroundings?''. Both are real questions with different answers, and
\emph{local RX}\index{local RX} answers the second by estimating a separate mean and covariance for
every pixel from the pixels around it.

\emph{The covariance is nearly singular.} On 165 correlated bands the smallest eigenvalues are
estimation noise, and inverting turns them into the largest weights (\cref{sec:rx:singular}). Diagonal
loading keeps the arithmetic from failing but cannot tell real low-variance directions from noise
(\cref{sec:rx:regularisation}). The repair the product ships is to reduce the dimension first, keeping
only directions that are well estimated and carry signal, and to run RX there. The obvious way to
reduce dimension is PCA; the right way, for a sensor whose noisiest bands are also its most variable,
is the \emph{maximum noise fraction} transform, MNF.

\begin{analogy}
A security guard deciding whether a car is out of place could compare it with every car in the city
(global) or with the cars parked on the same street (local). A sports car is unremarkable in the
city-wide statistics and alarming on a street of delivery vans. But the guard must not include the
suspicious car's own parking space, or the spaces right next to it where its passengers' other cars
may be, in ``the street'': those would make the street look like the suspect. That margin left out of
the comparison is the guard window of this chapter.
\end{analogy}

\section{The dual window}\label{sec:localrx:window}

For each pixel under test, local RX takes its background from a square ring: a square \emph{outer
window} centred on the pixel, with a smaller square \emph{inner window} (the \emph{guard
window}\index{guard window}) removed from its middle. This shape is called a \emph{dual
window}\index{dual window} or \emph{annulus}\index{annulus}.

\begin{definition}[Ring of half-sizes $o$ and $i$]
With outer half-size $o$ and inner half-size $i < o$, the outer window is the
$(2o + 1)\times(2o + 1)$ square of pixels within $o$ rows and $o$ columns of the test pixel, the guard
is the $(2i + 1)\times(2i + 1)$ square within $i$, and the background ring is the outer window minus
the guard.
\end{definition}

\begin{figure}[htbp]
\centering
\begin{tikzpicture}[scale=0.42]
  \foreach \x in {-4,...,4} \foreach \y in {-4,...,4} {
    \fill[spat!25] (\x-0.5,\y-0.5) rectangle (\x+0.5,\y+0.5);
  }
  \foreach \x in {-1,...,1} \foreach \y in {-1,...,1} {
    \fill[maskgrey!45] (\x-0.5,\y-0.5) rectangle (\x+0.5,\y+0.5);
  }
  \fill[loss] (-0.5,-0.5) rectangle (0.5,0.5);
  \draw[step=1,gray!60,shift={(0.5,0.5)}] (-5,-5) grid (4,4);
  \draw[very thick,spat] (-4.5,-4.5) rectangle (4.5,4.5);
  \draw[very thick,maskgrey!90!black] (-1.5,-1.5) rectangle (1.5,1.5);
  \draw[decorate,decoration={brace,amplitude=4pt}] (0.5,4.8) -- (4.5,4.8)
    node[midway,above=3pt,font=\scriptsize] {$o = 4$};
  \draw[decorate,decoration={brace,amplitude=3pt,mirror}] (0.5,-1.8) -- (1.5,-1.8)
    node[midway,below=3pt,font=\scriptsize] {$i = 1$};
  \node[right,font=\small,align=left] at (5.5,2) {\textcolor{loss}{\rule{6pt}{6pt}} pixel under test};
  \node[right,font=\small,align=left] at (5.5,0) {\textcolor{maskgrey}{\rule{6pt}{6pt}} guard, excluded: $3\times3 = 9$};
  \node[right,font=\small,align=left] at (5.5,-2) {\textcolor{spat!40}{\rule{6pt}{6pt}} background ring: $81 - 9 = 72$};
\end{tikzpicture}
\caption{A dual window with outer half-size $o = 4$ and inner half-size $i = 1$. The background is the
$9\times9$ outer square minus the $3\times3$ guard: $81 - 9 = 72$ pixels. The product's defaults,
$o = 25$ and $i = 5$, give a $51\times51$ square and an $11\times11$ guard.}
\label{fig:localrx:window}
\end{figure}

\subsection{Counting the ring}

The number of background pixels is the outer area minus the guard area. Factor it with the
difference of two squares, $a^2 - b^2 = (a - b)(a + b)$:
\begin{align}
n_{\text{ring}} &= (2o + 1)^2 - (2i + 1)^2 \notag\\
&= \big((2o + 1) - (2i + 1)\big)\big((2o + 1) + (2i + 1)\big) \notag\\
&= (2o - 2i)(2o + 2i + 2) \notag\\
&= 4\,(o - i)(o + i + 1). \label{eq:localrx:ring}
\end{align}

\begin{worked}[label={we:localrx:count}]{The default ring, and the ring at an edge}
\code{LocalRXDetector} defaults to \code{outer_window = 25} and \code{inner_window = 5}, which its
docstring defines as half-sizes (\code{app/detectors/local_rx_detector.py}, lines 53--54 and 155--156).
Directly:
\begin{align*}
\text{outer square} &= (2\times25 + 1)^2 = 51^2 = 2601\\
\text{guard} &= (2\times5 + 1)^2 = 11^2 = 121\\
n_{\text{ring}} &= 2601 - 121 = 2480 .
\end{align*}
By \cref{eq:localrx:ring}: $4\times(25 - 5)\times(25 + 5 + 1) = 4\times20\times31 = 2480$. The code
computes the same number as \code{max_bg = (2 * outer + 1) ** 2 - (2 * inner + 1) ** 2} (line 240).

Near the edge of the image the windows are cut off. At a corner pixel the outer window keeps only the
rows and columns inside the image, $0$ to $25$, so it is $26\times26$, and the guard is $6\times6$:
\[
n_{\text{corner}} = 26\times26 - 6\times6 = 676 - 36 = 640 .
\]
On the middle of the top edge the outer window is $26\times51$ and the guard $6\times11$:
$1326 - 66 = 1260$. A corner pixel has about a quarter of the full background, and pixels beside
cloud or nodata lose more.
\end{worked}

\section{Why the guard window exists}\label{sec:localrx:guard}

The guard is the subtle part, and leaving it out quietly ruins the detector. A real target is usually
larger than one pixel, and the sensor blurs it further into its neighbours. Without a guard, those
neighbours are \emph{also} target, and they go into the background statistics. The background mean
shifts towards the target, the background covariance inflates along exactly the direction that makes
the target unusual, and the Mahalanobis distance shrinks. The target partly whitens itself out of
detection.

\begin{worked}[label={we:localrx:guard}]{A $3\times3$ target with and without a guard}
One band, a $5\times5$ patch, outer half-size $o = 2$. The target is the $3\times3$ block of 25s in the
middle; the 16 pixels around it are background, eight at 20 and eight at 22:
\[
\begin{matrix}
20 & 22 & 20 & 22 & 20\\
22 & 25 & 25 & 25 & 22\\
20 & 25 & \mathbf{25} & 25 & 20\\
22 & 25 & 25 & 25 & 22\\
20 & 22 & 20 & 22 & 20
\end{matrix}
\]
We score the centre pixel (bold). Do both cases on paper.

\textbf{With a guard of half-size $i = 1$.} The guard is the whole $3\times3$ target, and the ring is
the 16 border pixels: $4\times(2 - 1)\times(2 + 1 + 1) = 16$. Their mean and variance:
\begin{align*}
\mu &= \frac{8\times20 + 8\times22}{16} = \frac{160 + 176}{16} = \frac{336}{16} = 21\\
\textstyle\sum(x - \mu)^2 &= 8\times(20 - 21)^2 + 8\times(22 - 21)^2 = 8 + 8 = 16\\
\sigma^2 &= \frac{16}{16 - 1} = \frac{16}{15} = 1.0667 .
\end{align*}
With one band RX is the squared z-score:
\[
\text{RX} = \frac{(25 - 21)^2}{16/15} = \frac{16\times15}{16} = 15 .
\]
That is far beyond the $\chi^2_1$ threshold for 0.1\% (10.83).

\textbf{Without a guard ($i = 0$, only the centre pixel excluded).} The ring is now 24 pixels: the 16
border pixels and the 8 other target pixels at 25.
\begin{align*}
\mu &= \frac{336 + 8\times25}{24} = \frac{336 + 200}{24} = \frac{536}{24} = 22.333\\
\textstyle\sum(x - \mu)^2 &= 8\times(20 - 22.333)^2 + 8\times(22 - 22.333)^2 + 8\times(25 - 22.333)^2\\
&= 8\times5.444 + 8\times0.111 + 8\times7.111 = 43.56 + 0.89 + 56.89 = 101.33\\
\sigma^2 &= \frac{101.33}{23} = 4.406\\
\text{RX} &= \frac{(25 - 22.333)^2}{4.406} = \frac{7.111}{4.406} = 1.61 .
\end{align*}
The same pixel drops from 15 to 1.61, an ordinary value. The target's own neighbours raised the
background mean from 21 to 22.3 and quadrupled its variance. The repository's
\code{batch_mahalanobis} returns 15.0000 and 1.6140 on this patch\tagC{} (with the default ridge of
$10^{-4}$, 14.9986 and 1.6140).
\end{worked}

\begin{keypoint}
Whatever the background window contains is, by definition, normal; the guard keeps the target out of
its own definition of normal.
\end{keypoint}

The guard is a compromise, not a guarantee. A target larger than the guard contaminates its ring
anyway; a guard so large that it holds every possible target leaves a ring far from the pixel, which
is a worse description of its surroundings. A target that fills the whole outer window becomes its own
background and vanishes. That is one reason to run global and local detectors side by side: they fail
on different targets.

\section{The local RX algorithm}\label{sec:localrx:algorithm}

For every valid pixel $\vect{x}$ at position $(r, c)$, \code{LocalRXDetector.detect} does the following.
\begin{enumerate}
  \item Collect the valid pixels of the ring around $(r, c)$: $\vect{b}_1, \dots, \vect{b}_n$.
  \item If $n$ is below \code{min_bg_pixels}, leave the pixel unscored (NaN).
  \item Compute the ring mean $\hat\muv = \frac1n\sum_j\vect{b}_j$ and covariance
    $\hat\Cov = \frac{1}{n - 1}\sum_j(\vect{b}_j - \hat\muv)(\vect{b}_j - \hat\muv)\T$.
  \item Score
  \begin{equation}\label{eq:localrx:lrx}
  \text{LRX}(\vect{x}) = (\vect{x} - \hat\muv)\T\,(\hat\Cov + \epsilon\mat{I})\inv\,(\vect{x} - \hat\muv),
  \qquad \epsilon = 10^{-4}.
  \end{equation}
\end{enumerate}
The source docstring states the same formula, ``Score = (x - $\mu$)$^\top$ ($\Sigma$ + $\lambda$I)$^{-1}$ (x -
$\mu$)'' (\code{local_rx_detector.py}, line 10), with the ridge called $\lambda$ there; this book calls it
$\epsilon$ to keep $\lambda$ for eigenvalues.

\subsection{Why at least $B + 1$ background pixels}

The default \code{min_bg_pixels} is the number of bands plus one, $B + 1$ (\code{local_rx_detector.py},
line 187; $10 + 1 = 11$ for the 10-component MNF version). The reason is a fact about sample
covariances.

\paragraph{Claim.} A covariance estimated from $n$ points has rank at most $n - 1$: at most $n - 1$ of
\index{rank}\index{covariance matrix!rank of a sample}
its eigenvalues are non-zero.

\noindent\textit{Proof.} By \cref{eq:linalg:covmat}, $\hat\Cov$ is a sum of $n$ outer products
$\vect{d}_j\vect{d}_j\T$. For any vector $\vect{v}$, $\hat\Cov\vect{v} = \frac{1}{n-1}\sum_j
\vect{d}_j(\vect{d}_j\T\vect{v})$ is a combination of the deviations $\vect{d}_j$, so every output of
$\hat\Cov$ lies in the space the deviations span. The deviations add to zero
(\cref{sec:linalg:covariance}), so the last one is minus the sum of the others, and they span at most
$n - 1$ dimensions. A matrix whose outputs fill at most $n - 1$ dimensions has at least $B - (n - 1)$
zero eigenvalues. \hfill$\square$

So with $n \le B$ background pixels the covariance is exactly singular and has no inverse; $n = B + 1$
is the least that can give a full-rank matrix. The script confirms it on random 10-dimensional points:
2 points give rank 1, 5 give rank 4, 10 give rank 9, 11 give rank 10\tagC. Full rank is necessary, not
sufficient: a covariance from 11 points in 10 dimensions is invertible but its small eigenvalues are
very poorly estimated, which is where the ridge earns its place. On 177 raw PRISMA bands the corner
pixel's 640 background pixels are enough for full rank but not for a trustworthy inverse.

\subsection{Batching, padding and the masked mean}
\index{masked mean}\index{local RX!batching}

Local RX needs one covariance, one solve and one quadratic form \emph{per pixel}. Doing this one pixel
at a time in Python is hopeless (\cref{sec:localrx:cost}), so the detector stacks a batch of pixels'
rings into one array of shape \shape{(N, max\_bg, B)} and computes all their statistics in a few tensor
operations (\code{app/detectors/_local_background.py}, \code{batch_mahalanobis}). Arrays must be
rectangular, but rings are not all the same size, so each row is padded to the full ring size
\code{max_bg} and a count array \code{n} says how many entries are real. A boolean mask
\code{arange(max_bg) < n[:, None]} marks the real entries.

\begin{worked}[label={we:localrx:padding}]{The masked mean must divide by the real count}
Two pixels, one band, \code{max_bg = 4}. Pixel 0 has four real background values; pixel 1 has two,
0.30 and 0.34, padded with two zeros.
\begin{align*}
\text{mask, pixel 0:}&\quad [0,1,2,3] < 4 \;\to\; [\text{T},\text{T},\text{T},\text{T}]\\
\text{mask, pixel 1:}&\quad [0,1,2,3] < 2 \;\to\; [\text{T},\text{T},\text{F},\text{F}]
\end{align*}
The correct mean of pixel 1 divides the masked sum by its real count, 2:
\[
\frac{0.30 + 0.34 + 0 + 0}{2} = \frac{0.64}{2} = 0.32 .
\]
Dividing by \code{max_bg} instead would give $0.64/4 = 0.16$, half the true value, because two of the
four slots are padding. The code divides by \code{mask.sum(dim=1)} (line 60--61), and it multiplies the
centred values by the mask again (line 64) so that padding contributes nothing to the covariance
either. With the correct statistics, a test value of 0.40 at pixel 1 scores
\[
\frac{(0.40 - 0.32)^2}{\sigma^2 + \epsilon},\quad
\sigma^2 = \frac{(0.30 - 0.32)^2 + (0.34 - 0.32)^2}{2 - 1} = \frac{0.0004 + 0.0004}{1} = 0.0008,
\]
so $0.0064/(0.0008 + 0.0001) = 0.0064/0.0009 = 7.11$, which is what \code{batch_mahalanobis}
returns\tagC. Edge pixels are exactly the ones with padding, and exactly where you least want
spurious detections.
\end{worked}

\subsection{Three practical details}

\paragraph{Precision depends on the device.} \code{batch_mahalanobis} computes in \code{float32} on a
GPU and \code{float64} on the CPU (line 47). The same scene can therefore give scores differing in the
last digits on different hardware, far below any threshold anyone sets, but enough to break a test
that demands exact equality. Test with a tolerance, or test the ranking.

\paragraph{Two masks.} The result object carries a \code{spatial_mask} (pixels \emph{eligible} for
\index{computed mask}
scoring: valid and kept) and a \code{computed_mask} (pixels that actually \emph{got} a score,
\code{local_rx_detector.py} line 368). They differ wherever a ring held fewer than \code{min_bg_pixels}
valid neighbours. Always reduce statistics over \code{computed_mask}. Treating an unscored pixel as
score zero biases every statistic, and since unscored pixels cluster at edges and beside cloud, the
bias is systematic rather than random.

\paragraph{Stride.} A \code{stride} greater than 1 scores only every \code{stride}-th row and column and
fills the rest by bilinear interpolation (\code{_interpolate_strided}). A stride of 2 scores a quarter
of the pixels, about four times faster, at the price of smoothing small targets.

\section{The cost of local RX}\label{sec:localrx:cost}
\index{local RX!cost}

\begin{worked}[label={we:localrx:cost}]{How much arithmetic one PRISMA scene costs}
A PRISMA scene is about $1210\times1219$ pixels\tagM{docs/lld/local-background.md}:
\[
1210\times1219 = 1{,}474{,}990 \text{ pixels}.
\]
For each, building a $B\times B$ covariance from 2480 ring pixels takes about $2480\times B^2$
multiply-adds (each ring pixel adds its $B^2$-entry outer product), and solving one $B\times B$ system
about $B^3/3$. Two cases:
\begin{center}
\begin{tabular}{rrrr}
\toprule
$B$ & covariance per pixel & solve per pixel & whole scene\\
\midrule
10  & $2480\times100 = 248{,}000$          & $\approx 333$       & $\approx 3.7\times10^{11}$\\
177 & $2480\times31{,}329 = 77{,}695{,}920$ & $\approx 1.85\times10^{6}$ & $\approx 1.2\times10^{14}$\\
\bottomrule
\end{tabular}
\end{center}
For comparison, global RX at $B = 10$ builds one covariance from all the pixels:
$1{,}474{,}990\times100 \approx 1.5\times10^{8}$ multiply-adds, more than two thousand times less than
local RX at the same $B$. Running local RX on 177 raw bands is about 300 times more expensive again than
running it in 10 MNF components\tagC.

Memory matters as much as arithmetic. The detector sizes its GPU batch from
\code{bytes_per_elem = (2 * max_bg * B + B * B) * 4} (\code{local_rx_detector.py}, lines 198--199): for
$B = 10$ that is $(2\times2480\times10 + 100)\times4 = 198{,}800$ bytes per pixel in the batch, and for
$B = 177$ it is $3{,}636{,}996$ bytes, so a batch that fits a GPU at $B = 10$ must be about 18 times
smaller at $B = 177$.
\end{worked}

The two repairs of this chapter therefore support each other. Reducing the dimension first (MNF) makes
local RX affordable, and makes each local covariance well estimated from 2480 pixels.

\section{Choosing a better space: why not PCA}\label{sec:localrx:whynotpca}

The resolver's comment says the fix for the near-singular covariance is to run RX ``in a 10-dim
noise-whitened subspace''. Why noise-whitened, and not the obvious choice, the top 10 principal
components of \cref{sec:linalg:pca}?

PCA keeps the directions of greatest variance. Variance is signal plus noise, and on an imaging
spectrometer the noise is not the same in every band: bands at the ends of the spectrum and in
absorption regions have little light and a lot of noise. If a noisy band's noise variance is large,
PCA ranks it high and keeps it, while a clean band with real but modest variation is ranked low and
thrown away. What you want to rank by is not variance but \emph{signal-to-noise ratio}.

\subsection{Signal plus noise}

Model every pixel as $\vect{x} = \vect{s} + \vect{n}$: a signal $\vect{s}$ (what the ground looks like)
plus sensor noise $\vect{n}$, with the noise independent of the signal and of mean zero. The covariance
of a sum of independent parts is the sum of their covariances. (Expand
$(\vect{d}_s + \vect{d}_n)(\vect{d}_s + \vect{d}_n)\T = \vect{d}_s\vect{d}_s\T + \vect{d}_s\vect{d}_n\T +
\vect{d}_n\vect{d}_s\T + \vect{d}_n\vect{d}_n\T$ and average: the cross terms average to zero for
independent parts.) So
\begin{equation}\label{eq:localrx:sigsum}
\Cov = \Cov_s + \Cov_n .
\end{equation}
Along a unit direction $\vect{a}$, by \cref{eq:linalg:dirvar}, the signal variance is
$\vect{a}\T\Cov_s\vect{a}$ and the noise variance $\vect{a}\T\Cov_n\vect{a}$.

\begin{definition}[Noise fraction and signal-to-noise ratio]\index{noise fraction}\index{signal-to-noise ratio}
The \emph{noise fraction} of a direction $\vect{a}$ is the share of its variance that is noise,
$\vect{a}\T\Cov_n\vect{a} / \vect{a}\T\Cov\vect{a}$. Its \emph{signal-to-noise ratio} (SNR) is
$\vect{a}\T\Cov_s\vect{a} / \vect{a}\T\Cov_n\vect{a}$.
\end{definition}

The \emph{maximum noise fraction}\index{maximum noise fraction} transform of \citet{green1988transformation}
finds the directions in order of noise fraction; ordered from least noisy to most, they are the
directions of highest SNR first. The same transform is often called the \emph{minimum noise fraction}
transform, which is the name in the course and in the repository's docstrings; both names refer to the
same directions, listed from opposite ends.

\section{Deriving MNF}\label{sec:localrx:mnf}

\subsection{The generalised eigenproblem}

Maximise the ratio of total variance to noise variance along $\vect{a}$:
\begin{equation}\label{eq:localrx:rayleigh}
\max_{\vect{a}}\ R(\vect{a}) = \frac{\vect{a}\T\Cov\vect{a}}{\vect{a}\T\Cov_n\vect{a}} .
\end{equation}
Because $\Cov = \Cov_s + \Cov_n$, $R(\vect{a}) = 1 + \text{SNR}(\vect{a})$, so maximising $R$ maximises SNR
and minimises the noise fraction $1/R$. Scaling $\vect{a}$ does not change $R$ (numerator and
denominator both scale by the square), so fix the scale by requiring $\vect{a}\T\Cov_n\vect{a} = 1$ and
maximise the numerator, exactly as PCA fixed $\vect{u}\T\vect{u} = 1$ in \cref{sec:linalg:pca}. The
Lagrangian is
\[
L(\vect{a}, \lambda) = \vect{a}\T\Cov\vect{a} - \lambda(\vect{a}\T\Cov_n\vect{a} - 1).
\]
The gradient of a quadratic form $\vect{a}\T\mat{S}\vect{a}$ with symmetric $\mat{S}$ is $2\mat{S}\vect{a}$
(\cref{sec:linalg:pca}), so setting the gradient of $L$ to zero gives
\begin{align}
2\Cov\vect{a} - 2\lambda\Cov_n\vect{a} &= \vect{0} \notag\\
\Cov\vect{a} &= \lambda\,\Cov_n\vect{a} . \label{eq:localrx:geneig}
\end{align}
This is a \emph{generalised eigenproblem}\index{generalised eigenproblem}: $\lambda$ and $\vect{a}$ are
a generalised eigenvalue and eigenvector of the pair $(\Cov, \Cov_n)$. If $\Cov_n$ is invertible,
multiplying both sides on the left by $\Cov_n\inv$ gives the ordinary eigenproblem
$\Cov_n\inv\Cov\,\vect{a} = \lambda\vect{a}$. Multiplying \cref{eq:localrx:geneig} on the left by
$\vect{a}\T$ shows what $\lambda$ means:
$\vect{a}\T\Cov\vect{a} = \lambda\,\vect{a}\T\Cov_n\vect{a}$, so $\lambda = R(\vect{a}) = 1 +
\text{SNR}(\vect{a})$. The best direction is the generalised eigenvector with the largest $\lambda$.

\subsection{Solving it: noise-whiten, then PCA}

$\Cov_n\inv\Cov$ is not symmetric, so the spectral theorem does not apply to it directly. The trick is
to change variables so that it does. Let $\mat{W}_n = \Cov_n^{-1/2}$ be a whitening matrix for the
\emph{noise} (\cref{sec:linalg:whitening}, symmetric form), so that $\mat{W}_n\Cov_n\mat{W}_n = \mat{I}$.
Write $\vect{a} = \mat{W}_n\vect{v}$ and substitute into \cref{eq:localrx:geneig}:
\begin{align*}
\Cov\mat{W}_n\vect{v} &= \lambda\,\Cov_n\mat{W}_n\vect{v}\\
\mat{W}_n\Cov\mat{W}_n\vect{v} &= \lambda\,\mat{W}_n\Cov_n\mat{W}_n\vect{v}
  && \text{multiply on the left by } \mat{W}_n\\
\left(\mat{W}_n\Cov\mat{W}_n\right)\vect{v} &= \lambda\,\vect{v}
  && \mat{W}_n\Cov_n\mat{W}_n = \mat{I}.
\end{align*}
The matrix $\mat{W}_n\Cov\mat{W}_n$ is symmetric: it is the covariance of the noise-whitened data
$\mat{W}_n\vect{x}$ (\cref{eq:linalg:covtransform}). So MNF is two steps you already know:
\begin{enumerate}
  \item \textbf{Noise-whiten}: transform the data so that the noise has identity covariance, equal in
    every direction.
  \item \textbf{PCA} in the whitened space. Now that every direction has noise variance 1, the
    direction of greatest variance is the direction of greatest $1 + \text{SNR}$.
\end{enumerate}
The MNF directions in the original band space are $\vect{a}_k = \mat{W}_n\vect{v}_k$, and the $k$-th
\emph{MNF component} of a pixel is $y_k = \vect{a}_k\T(\vect{x} - \muv) = \vect{v}_k\T\mat{W}_n(\vect{x} -
\muv)$. This is exactly the code:

\begin{lstlisting}
# app/detectors/mnf_compression_detector.py, _compute_mnf_transform (lines 315-344)
centered = pixels - self._mnf_mean                     # (N, B)
noise_eigvals, noise_eigvecs = np.linalg.eigh(noise_cov)
noise_eigvals = np.maximum(noise_eigvals, 1e-10)       # floor the noise eigenvalues
# W_n, (B, B):
whiten = noise_eigvecs @ np.diag(1.0 / np.sqrt(noise_eigvals)) @ noise_eigvecs.T
whitened = centered @ whiten.T                         # step 1, (N, B)
# step 2, PCA of the noise-whitened data:
eigvals, eigvecs = np.linalg.eigh(np.cov(whitened, rowvar=False))
order = np.argsort(eigvals)[::-1]                      # largest 1 + SNR first
combined = (eigvecs[:, :n_components].T) @ whiten      # rows a_k^T, (k, B)
\end{lstlisting}

The symmetric form $\mat{V}\mat{\Lambda}^{-1/2}\mat{V}\T$ used here is one of many noise-whitening
matrices; any other is a rotation of it, $\mat{R}\mat{W}_n$ with $\mat{R}$ orthogonal. A rotation
before PCA rotates the eigenvectors $\vect{v}_k$ by the same $\mat{R}$ and leaves the eigenvalues
unchanged, so the final directions $\vect{a}_k$ and components $y_k$ are the same whichever whitening you
pick.

\begin{worked}[label={we:localrx:mnf2}]{MNF against PCA on two bands}
Band 1 is variable but noisy; band 2 is quiet but clean. Suppose the total and noise covariances are
\[
\Cov = \begin{pmatrix} 8 & 2\\ 2 & 1.25\end{pmatrix},\qquad
\Cov_n = \begin{pmatrix} 4 & 0\\ 0 & 0.25\end{pmatrix},
\qquad\text{so}\qquad
\Cov_s = \Cov - \Cov_n = \begin{pmatrix} 4 & 2\\ 2 & 1\end{pmatrix}.
\]
Half of band 1's variance is noise (4 of 8); a fifth of band 2's is (0.25 of 1.25). The signal
covariance has determinant $4\times1 - 2\times2 = 0$: all the signal lies along one direction,
$(2, 1)$, as if one material's abundance varied across the scene. Do this on paper.

\textbf{Step 1, noise-whiten.} $\Cov_n$ is diagonal, so
$\mat{W}_n = \Cov_n^{-1/2} = \diag(1/\sqrt4, 1/\sqrt{0.25}) = \diag(0.5, 2)$. The whitened covariance is
$\mat{W}_n\Cov\mat{W}_n$; entry $(i,j)$ is $\Sigma_{ij}$ times the $i$-th and $j$-th diagonal entries of
$\mat{W}_n$:
\begin{align*}
(1,1):&\quad 8\times0.5\times0.5 = 2\\
(1,2):&\quad 2\times0.5\times2 = 2\\
(2,2):&\quad 1.25\times2\times2 = 5
\end{align*}
so $\mat{W}_n\Cov\mat{W}_n = \begin{psmallmatrix} 2 & 2\\ 2 & 5\end{psmallmatrix}$.

\textbf{Step 2, PCA.} Trace $2 + 5 = 7$, determinant $2\times5 - 2\times2 = 6$, characteristic
polynomial $\lambda^2 - 7\lambda + 6 = (\lambda - 6)(\lambda - 1)$: $\lambda_1 = 6$, $\lambda_2 = 1$. For
$\lambda_1 = 6$: $(2 - 6)v_1 + 2v_2 = 0$, so $v_2 = 2v_1$ and $\vect{v}_1 = (1, 2)/\sqrt5$. For
$\lambda_2 = 1$: $(2 - 1)v_1 + 2v_2 = 0$, so $\vect{v}_2 = (2, -1)/\sqrt5$.

\textbf{Step 3, back to bands.} $\vect{a}_1 = \mat{W}_n\vect{v}_1 = (0.5\times1, 2\times2)/\sqrt5 = (0.5,
4)/\sqrt5 = (0.224, 1.789)$. The first MNF direction weights the clean band 8 times as heavily as the
noisy one.

\textbf{Step 4, read the eigenvalues.} $\lambda_1 = 6 = 1 + \text{SNR}$, so the first component has
SNR 5. Check directly: $\Cov_s\vect{a}_1 = (4\times0.5 + 2\times4,\ 2\times0.5 + 1\times4)/\sqrt5 =
(10, 5)/\sqrt5$, so the signal variance is $\vect{a}_1\T\Cov_s\vect{a}_1 = (0.5\times10 + 4\times5)/5 =
25/5 = 5$; the noise variance is $(4\times0.25 + 0.25\times16)/5 = (1 + 4)/5 = 1$; SNR $= 5/1 = 5$.
The second component has $\lambda_2 = 1$: SNR 0, pure noise, as it must be when all the signal lies
along one direction.

\textbf{Compare PCA on the raw bands.} $\Cov$ has trace $9.25$ and determinant
$8\times1.25 - 4 = 6$, so its eigenvalues are $(9.25 \pm \sqrt{9.25^2 - 24})/2 = (9.25 \pm 7.846)/2$:
8.548 and 0.702. The first principal direction is $(0.964, 0.264)$, mostly the noisy band 1. Its signal
variance is 4.81 and its noise variance 3.74\tagC: SNR 1.29. PCA's best direction is a quarter as clean
as MNF's.
\end{worked}

\subsection{Estimating the noise from neighbouring pixels}
\index{noise covariance}\index{maximum noise fraction!noise estimate}

MNF needs $\Cov_n$, and nobody hands you the sensor noise. The \emph{shift-difference}
method\index{shift difference} of \citet{green1988transformation} estimates it from the image itself.
Adjacent pixels usually show nearly the same ground, so their \emph{difference} is mostly noise. Write
two horizontal neighbours as $\vect{x}(p) = \vect{s}(p) + \vect{n}(p)$ and $\vect{x}(p') = \vect{s}(p') +
\vect{n}(p')$. If $\vect{s}(p') \approx \vect{s}(p)$,
\[
\vect{x}(p') - \vect{x}(p) \approx \vect{n}(p') - \vect{n}(p).
\]
If the noise at the two pixels is independent with the same covariance $\Cov_n$, the covariance of
the difference is
\begin{align*}
\Cov\big(\vect{n}(p') - \vect{n}(p)\big)
&= \Cov(\vect{n}(p')) + \Cov(\vect{n}(p)) - (\text{cross terms, zero for independent noise})\\
&= \Cov_n + \Cov_n = 2\,\Cov_n ,
\end{align*}
using \cref{eq:linalg:covtransform} with $\mat{A} = (\mat{I}\ \ {-\mat{I}})$ applied to the stacked pair.
So $\Cov_n \approx \frac12\Cov(\text{differences})$: the line
\code{noise_cov = np.cov(diff_pixels, rowvar=False) / 2.0}
(\code{mnf_compression_detector.py}, line 291), with differences taken between each valid pixel and its
right-hand neighbour (lines 283--285).

\begin{worked}[label={we:localrx:noise}]{A noise covariance from a strip of five pixels}
Two bands, one row of five pixels:
\begin{center}
\begin{tabular}{cccccc}
\toprule
& $p_1$ & $p_2$ & $p_3$ & $p_4$ & $p_5$\\
\midrule
band 1 & 10.0 & 10.2 & 9.9  & 10.1 & 10.0\\
band 2 & 20.0 & 20.1 & 20.1 & 20.0 & 20.0\\
\bottomrule
\end{tabular}
\end{center}
Horizontal differences, right minus left:
\begin{align*}
\text{band 1:}&\quad 0.2,\ -0.3,\ 0.2,\ -0.1 \qquad(\text{sum } 0)\\
\text{band 2:}&\quad 0.1,\ 0.0,\ -0.1,\ 0.0 \qquad(\text{sum } 0).
\end{align*}
Their means are 0, so the covariance of the four difference vectors (divisor $4 - 1 = 3$) is
\begin{align*}
\text{var}_1 &= \tfrac13(0.04 + 0.09 + 0.04 + 0.01) = \tfrac13\times0.18 = 0.06\\
\text{var}_2 &= \tfrac13(0.01 + 0 + 0.01 + 0) = \tfrac13\times0.02 = 0.00667\\
\text{cov}_{12} &= \tfrac13\big(0.2\times0.1 + (-0.3)\times0 + 0.2\times(-0.1) + (-0.1)\times0\big)
  = \tfrac13(0.02 - 0.02) = 0 .
\end{align*}
Halving: $\hat\Cov_n = \diag(0.03, 0.00333)$, noise standard deviations $\sqrt{0.03} = 0.173$ and
$\sqrt{0.00333} = 0.058$. Band 1 is three times as noisy. The repository's
\code{_estimate_noise_covariance}, called on this strip, returns the same matrix\tagC.
\end{worked}

The assumption behind the estimate fails wherever neighbours really differ: along a coastline, a road
edge, a field boundary. Those differences are counted as noise, inflating $\hat\Cov_n$ in the directions
of the edges' contrast. Over a whole scene, edges are a small fraction of neighbour pairs and the
estimate is useful; and it needs no calibration data. One consequence is worth noticing: an
imaging spectrometer such as PRISMA records each image column with its own detector element, so a
column-to-column calibration offset (a stripe) appears in exactly these horizontal differences. The
estimate counts striping as noise, which is why the MNF local detector's docstring credits it with
``suppressing sensor noise and striping artefacts'' (\code{mnf_compression_lrx_detector.py}, lines
6--8).

\subsection{The eigenvalue diagnostic}
\index{maximum noise fraction!eigenvalues}

Because each MNF eigenvalue is $1 + \text{SNR}$ of its component, the list of eigenvalues says how much
signal the scene has and in how many directions. Signal components have eigenvalues well above 1; pure
noise components sit near 1. The anomaly course states the practical rule: the eigenvalues should
\emph{fall away sharply}, a few components carrying the structure; nearly flat eigenvalues mean the
compression found no dominant signal and the scores deserve suspicion. That is why
\code{MNFCompressionRXResult} carries \code{mnf_eigenvalues} and why the detector logs the ratio of
the first to the last retained one (\code{mnf_compression_detector.py}, lines 346--349).

\begin{worked}[label={we:localrx:spectrum}]{Eigenvalues of PCA and MNF on a synthetic cube}
The script builds a $12$-band, $50\times50$ synthetic cube from three smooth spatial patterns mixed by
three random spectra, plus independent noise whose standard deviation rises from 0.05 in band 1 to 0.6
in band 12, and runs the repository's own MNF code on it\tagC. The noise estimate is close: for band 12
the true noise variance is $0.6^2 = 0.36$ and the estimate is 0.350; for band 4, 0.040 against
0.042. The two eigenvalue lists:
\begin{center}
\begin{tabular}{lp{0.72\linewidth}}
\toprule
PCA of the cube & 1.26, 0.53, 0.31, 0.26, 0.20, 0.17, 0.14, 0.085, 0.059, 0.035, 0.024, 0.010\\
MNF & 65.4, 4.02, 2.18, \quad 1.07, 1.06, 1.02, 1.00, 0.99, 0.98, 0.97, 0.95, 0.93\\
\bottomrule
\end{tabular}
\end{center}
PCA's list falls smoothly: you cannot see that the cube has exactly three signals, because the noisy
bands supply variance in every direction. MNF's list has three eigenvalues well above 1 and nine
within 7\% of 1: three signal components and nine noise components, read straight off the table.
\end{worked}

\section{MNF-RX}\label{sec:localrx:mnfrx}

\emph{MNF-RX}\index{MNF-RX} projects every valid pixel onto the first $k$ MNF directions and runs global
RX on the $k$ components. The product's default is $k = 10$
(\code{DEFAULT_N_COMPONENTS = 10}, \code{mnf_compression_detector.py} line 33; the worker calls
\code{fit()} with no arguments). To see what this buys, you need one more fact about the Mahalanobis
distance.

\subsection{Mahalanobis distance ignores invertible changes of coordinates}
\index{Mahalanobis distance!invariance}

\paragraph{Claim.} Let $\vect{y} = \mat{A}\vect{x}$ with $\mat{A}$ square and invertible. Then the
Mahalanobis distance of $\vect{y}$ in its own statistics equals that of $\vect{x}$ in its own.

\noindent\textit{Proof.} The mean of $\vect{y}$ is $\mat{A}\muv$ and its covariance
$\mat{A}\Cov\mat{A}\T$ (\cref{eq:linalg:covtransform}). The inverse of a product of invertible matrices
is the product of inverses in reverse order, $(\mat{A}\Cov\mat{A}\T)\inv = (\mat{A}\T)\inv\Cov\inv\mat{A}\inv$
(check: multiplying by $\mat{A}\Cov\mat{A}\T$ cancels the factors pair by pair from the middle). With
$\vect{d} = \vect{x} - \muv$, use the transpose of a product and then the inverse of a product:
\begin{align*}
(\mat{A}\vect{d})\T(\mat{A}\Cov\mat{A}\T)\inv(\mat{A}\vect{d})
&= \vect{d}\T\mat{A}\T(\mat{A}\T)\inv\,\Cov\inv\,\mat{A}\inv\mat{A}\vect{d}\\
&= \vect{d}\T\,\mat{I}\,\Cov\inv\,\mat{I}\,\vect{d} = \vect{d}\T\Cov\inv\vect{d}.
\end{align*}
\hfill$\square$

The MNF transform with all $B$ components kept is an invertible matrix, so \textbf{MNF-RX with every
component kept is exactly plain RX}. The script confirms it on the synthetic cube: the two score maps
differ by at most $7\times10^{-14}$\tagC, rounding error. Every benefit of MNF-RX therefore comes from
\emph{discarding} components, not from the transform itself. MNF's job is to sort the directions so
that the ones you discard are the noise.

\subsection{What MNF-RX computes}

The MNF components are uncorrelated over the pixels used to fit them, with variances equal to the MNF
eigenvalues: their covariance is $\mat{V}_k\T\mat{W}_n\Cov\mat{W}_n\mat{V}_k = \diag(\lambda_1, \dots,
\lambda_k)$, by the same argument as \cref{we:linalg:pcacoords}. When \code{spectral.rx} re-estimates the
covariance of the $k$ components on those same pixels, it finds this diagonal matrix, so
\begin{equation}\label{eq:localrx:mnfrx}
\text{MNF-RX}(\vect{x}) = \sum_{j=1}^{k}\frac{y_j^2}{\lambda_j},
\qquad y_j = \vect{a}_j\T(\vect{x} - \muv).
\end{equation}
The script checks this on the synthetic cube with $k = 3$: the library's scores and the sum agree to
$5\times10^{-14}$, and the mean in-sample score is $2.9988 = 3\times(N-1)/N$\tagC, as
\cref{sec:rx:chi2} predicts.

Now the numerical point. The denominators in \cref{eq:localrx:mnfrx} are the $k$ largest MNF
eigenvalues, each $1 + \text{SNR} \ge$ roughly 1. No denominator can be tiny, so nothing explodes; the
condition number of the $10\times10$ covariance is $\lambda_1/\lambda_{10}$, the ratio the detector
logs, and with $\lambda_{10}$ near or above 1 it is bounded by the strongest component's SNR. Compare
plain RX, whose denominators include the smallest of 165 eigenvalues of the raw covariance. Under a
Gaussian background the scores follow $\chi^2_{10}$, mean 10 (\cref{sec:rx:chi2}).

There is a statistical point as well, and it explains why discarding helps even when nothing
explodes. A discarded MNF component with eigenvalue near 1 is almost pure noise; kept, it adds a
squared standard normal to every pixel's score and nothing to a target's. Under a Gaussian background a
score in $C$ dimensions has standard deviation $\sqrt{2C}$ (\cref{sec:rx:chi2}), and a target must stand
out against that spread. Suppose a target adds 30 to the score along one signal component. In $C = 10$
dimensions the background spread is $\sqrt{20} = 4.5$, and 30 is $30/4.5 = 6.7$ spreads; in $C = 165$
dimensions the spread is $\sqrt{330} = 18.2$, and the same 30 is $30/18.2 = 1.6$ spreads. Every noise
component kept dilutes the signal components.

\begin{sourcenote}[: MNF-RX in the product]
The worker builds the classical detector, calls \code{detector.fit()}, which estimates the MNF
transform from all spatially valid pixels, and only then narrows \code{detector._spatial_mask} to the
upstream keep mask (segmentation or cloud mask) before \code{detect()}
(\code{backend/allotrope/action_types/_anomaly_scoring_run.py}, lines 361--392). So the MNF directions
and the noise estimate come from all valid pixels, cloud included, while the RX mean and covariance in
MNF space come from the kept pixels only. In that case the components are no longer exactly
uncorrelated over the pixels RX sees, and \code{spectral.rx}'s full $10\times10$ inverse, not the
diagonal sum of \cref{eq:localrx:mnfrx}, is what is computed; the two agree only when no pixels are
masked. The course describes MNF-RX without this detail.
\end{sourcenote}

\subsection{MNF local RX}
\index{MNF local RX}

\code{MNFCompressionLRXDetector} combines both repairs: project onto the top $k = 10$ MNF components,
then run the local RX of \cref{sec:localrx:algorithm} on the $k$-band cube, with the same defaults
(outer half-size 25, guard half-size 5, ridge $10^{-4}$, \code{min_bg_pixels} $= k + 1 = 11$;
\code{mnf_compression_lrx_detector.py}, lines 36--39 and 215--219). Two consequences of the earlier
sections:
\begin{itemize}
  \item Each ring has 2480 pixels to estimate a $10\times10$ covariance with
    $10\times11/2 = 55$ distinct entries: about 45 pixels per number, well determined.
  \item The ridge $\epsilon = 10^{-4}$ is negligible here. The MNF components have variances of order
    1 or more, so adding $10^{-4}$ changes each eigenvalue by a fraction of $10^{-4}$ or less. In this
    space the ridge does what it should: it only matters when a ring is so small that its covariance
    is singular.
\end{itemize}
It is registered (\code{ADModel.MNF_LRX}) but, like plain local RX, it is not one of the product's
resolved classical models (\code{mnf_rx} and \code{thermal_grx}).

\section{PCA-RX, the evaluation baseline}\label{sec:localrx:pcarx}

One more Mahalanobis detector appears throughout Part VII as the classical yardstick. In the
evaluation framework (\code{app/evaluation/scorers.py}, lines 39--40) the scorer \code{pca_rx} is the
residual score of \cref{ch:residualmaha} applied to the cube with a reconstruction of zeros: the
``residual'' is then the cube itself.

\begin{definition}[PCA-RX]\index{PCA-RX}
Fit PCA to the valid pixels of the cube, keep the top $d = 30$ principal directions, and whiten them
(\cref{sec:linalg:whitening}): $\vect{z}(p) = \mat{\Lambda}_d^{-1/2}\mat{U}_d\T(\vect{x}(p) - \muv)$. The
score is the squared distance of a pixel's whitened components from the mean of the whitened
components over a ring around it:
\begin{equation}\label{eq:localrx:pcarx}
\text{PCA-RX}(p) = \norm{\vect{z}(p) - \bar{\vect{z}}_{\text{ring}}(p)}^2 .
\end{equation}
\end{definition}

It differs from \code{LocalRXDetector} in four ways, each worth knowing:
\begin{enumerate}
  \item \textbf{PCA, not MNF.} The 30 kept directions are those of greatest variance. Truncating at
    30 rather than 10 keeps more structure; the price is that noise-dominated high-variance
    directions, if any, are kept too.
  \item \textbf{Local mean, global covariance.} Only the mean comes from the ring; the whitening uses one
    covariance for the scene, fitted on at most 200,000 valid pixels drawn with seed 0
    (\code{residual_mahalanobis.py}, lines 35 and 38--46). A local mean costs a few sums per pixel
    instead of a covariance and a solve, and \cref{ch:residualmaha} computes it for every pixel at
    once with integral images.
  \item \textbf{Full window widths.} \code{LOCAL_WINDOW = (21, 61)} gives the guard and outer
    \emph{full} widths (line 33: ``inner (guard) and outer full widths''), not half-sizes. The ring has
    $61^2 - 21^2 = 3721 - 441 = 3280$ pixels. In \code{LocalRXDetector}'s half-size convention this would
    be $o = 30$, $i = 10$: $4\times20\times41 = 3280$.
  \item \textbf{A fallback instead of NaN.} A pixel whose ring holds fewer than $2d = 60$ valid pixels
    (\code{MIN_RING_FACTOR = 2}) is scored against the global mean, $\norm{\vect{z}(p)}^2$, and pixels
    outside the valid mask score 0, not NaN.
\end{enumerate}

\begin{worked}[label={we:localrx:pcarx}]{PCA-RX on a tiny whitened patch}
Take $d = 1$ and suppose the whitened component $z$ is 1 everywhere on a $7\times7$ patch except 4 at
the centre, with windows $(1, 5)$: a guard of full width 1 (the pixel itself) and an outer square of full
width 5. The ring around the centre has $5^2 - 1^2 = 24$ pixels, all with $z = 1$, so its mean is 1 and
\[
\text{PCA-RX}(\text{centre}) = (4 - 1)^2 = 9 .
\]
\code{local_mean_score} returns 9.0 on this patch\tagC. Note what is \emph{not} divided: there is no
ring variance in the denominator, because $z$ is already in units of the scene's global spread.
\end{worked}

\begin{drift}
PRISMA labelled bench, mean pixel ROC-AUC over four scenes, report generated 2026-10-15, production Indradhanu
v0.2.0 (\file{research/experiments/dual-window-mahalanobis/README.md}): MNF-RX with 10 components 0.752,
local RX with a dual window on the cube 0.743, PCA-RX with 30 components on the cube 0.831, residual
Mahalanobis on Indradhanu's residual 0.839 (global mean) and 0.854 (local mean, 21/61)\tagM{research/experiments/dual-window-mahalanobis/README.md}. On implant
bench-v1 (evaluation stage 3, 2026-10-17) the model-free PCA-RX caught 30.2\% of objects in the top 1\% against 26.9\% for FT03 with
residual Mahalanobis, while on the real objects its pAUC was 0.128 against FT03's
0.158\tagM{docs/training/anomaly-detection/19-running-an-evaluation.md}. \Cref{ch:benchmarks,ch:runeval}
explain what that reversal says about the bench.
\end{drift}

The README draws the lesson that matters for this chapter: PCA30-RX on the cube beats MNF-RX on both
datasets it was tried on, ``which suggests MNF-RX's 10 components are too few. That is a separate,
untested hypothesis.'' Read it as a caution, not a verdict: MNF-RX is the product's classical
hyperspectral detector because it is numerically safe, not because it was shown to be the best
classical detector.

\section{Global and local together}\label{sec:localrx:together}
\index{statistical ensembler}

Global and local detectors fail differently. Global RX misses a target that is ordinary for the scene
but strange for its street; local RX misses a target larger than its guard, and a target that sits in a
region it shares with nothing else. The repository has a detector that fuses the two,
\code{StatisticalEnsembler} (\code{app/detectors/statistical_ensembler.py}): it runs global and local RX,
turns each score map into ranks between 0 and 1 (its empirical cumulative distribution), and combines
them by product (the default: a pixel must look anomalous to both), maximum or mean. It is not
registered in the detector factory and is used only from notebooks. The same rank-fusion idea returns in
the scoring experiments of \cref{ch:scoringgains}.

\begin{confusions}
\confusion{Why not always use local RX, if it is more sensitive?}{It is far more expensive (one
covariance and solve per pixel), it misses targets larger than its guard, and it is blind to a target
that fills its ring. Global and local answer different questions.}
\confusion{Is the ring square or circular?}{Square: two nested squares, the inner one excluded. A
circle would be marginally more principled and slower to index.}
\confusion{Are outer\_window = 25 and LOCAL\_WINDOW = (21, 61) the same kind of number?}{No.
\texttt{LocalRXDetector}'s windows are half-sizes, so 25 means a $51\times51$ square; the residual score's
and PCA-RX's are full widths, so 61 means a $61\times61$ square. The \texttt{spectral} library's own
\texttt{rx(window=(inner, outer))} also takes full widths.}
\confusion{MNF is PCA, isn't it?}{PCA after noise-whitening. On the raw bands PCA maximises variance;
MNF maximises signal-to-noise. Where the noisiest bands are also the most variable, PCA keeps exactly the
wrong directions (\cref{we:localrx:mnf2}).}
\confusion{Does MNF-RX change the RX score of a pixel?}{Only by discarding components. With all
components kept it equals plain RX exactly; the Mahalanobis distance does not care about invertible
changes of coordinates.}
\confusion{Estimating noise from neighbour differences seems crude.}{It assumes adjacent pixels show the
same ground, which fails at edges, and counts column stripes as noise. It needs no calibration data and
is good enough over a whole scene, which is why it has been standard since 1988.}
\confusion{What is min\_bg\_pixels for?}{A covariance from $n$ points has rank at most $n - 1$, so with
$B$ bands you need at least $B + 1$ background pixels for an invertible covariance. Pixels with fewer
are left unscored (NaN), and that is why \texttt{computed\_mask} exists.}
\end{confusions}

\begin{checkyourself}
\item\label{ex:localrx:1} Compute the number of background pixels for half-sizes $o = 15$, $i = 3$, both
  directly and with \cref{eq:localrx:ring}. How many does a corner pixel have?
\item\label{ex:localrx:2} In the $5\times5$ patch of \cref{we:localrx:guard}, replace the eight target
  neighbours by background values (four 20s and four 22s) so that only the centre is 25. Score the
  centre with $i = 0$. What does this tell you about when a guard is needed?
\item\label{ex:localrx:3} A pixel has 3 real background values padded to 6: 0.20, 0.22, 0.24. Compute the
  correct mean and the mean you would get by dividing by 6, and the correct (unregularised) one-band
  score of a test value 0.30.
\item\label{ex:localrx:4} For $\Cov = \begin{psmallmatrix} 3 & 1\\ 1 & 1.5\end{psmallmatrix}$ and
  $\Cov_n = \begin{psmallmatrix} 2 & 0\\ 0 & 0.5\end{psmallmatrix}$, compute the noise-whitened covariance,
  its eigenvalues, and the SNR of the first MNF component. Which band does the first MNF direction weight
  more heavily?
\item\label{ex:localrx:5} Prove that MNF-RX with all components kept equals plain RX, starting from
  $\vect{y} = \mat{A}(\vect{x} - \muv)$ with $\mat{A}$ invertible. Why, then, does MNF-RX help?
\item\label{ex:localrx:6} How many distinct numbers does a $30\times30$ covariance hold? How many valid
  pixels per number does PCA-RX's 200,000-pixel sample give, and how many per number would a local
  $30$-dimensional covariance from a 3280-pixel ring give?
\end{checkyourself}

\nextchapter{ch:recon}
